\documentclass[12pt,a4paper]{article}

% ── Packages ──────────────────────────────────────────────────────────────────
\usepackage[utf8]{inputenc}
\usepackage[T1]{fontenc}
\usepackage[czech]{babel}
\usepackage[margin=2.5cm, headheight=14.5pt]{geometry}
\usepackage{amsmath}
\usepackage{graphicx}
\usepackage{booktabs}
\usepackage{longtable}
\usepackage{array}
\usepackage{enumitem}
\usepackage{titlesec}
\usepackage{fancyhdr}
\usepackage{xcolor}
\usepackage{listings}
\usepackage{parskip}
\usepackage{tikz}
\usetikzlibrary{
  shapes.geometric, shapes.multipart,
  arrows.meta, positioning, fit,
  calc, backgrounds,
  decorations.pathreplacing,
  matrix
}
\usepackage{hyperref}

% ── Colours ───────────────────────────────────────────────────────────────────
\definecolor{codeblue}{rgb}{0.13,0.29,0.53}
\definecolor{codegray}{rgb}{0.5,0.5,0.5}
\definecolor{codegreen}{rgb}{0.0,0.5,0.0}
\definecolor{backcolour}{rgb}{0.97,0.97,0.97}
\definecolor{teal}{RGB}{0,137,123}
\definecolor{tealdark}{RGB}{0,105,94}
\definecolor{tealpale}{RGB}{224,245,243}
\definecolor{layerA}{RGB}{230,242,255}
\definecolor{layerB}{RGB}{240,255,240}
\definecolor{layerC}{RGB}{255,255,225}
\definecolor{layerD}{RGB}{255,240,240}
\definecolor{boxborder}{RGB}{180,180,180}

% ── TikZ common styles ────────────────────────────────────────────────────────
\tikzset{
  layer/.style={
    rectangle, draw=boxborder, line width=0.8pt,
    rounded corners=4pt, align=center,
    minimum width=11cm, minimum height=1.4cm,
    text width=10.5cm, font=\small\bfseries
  },
  projectbox/.style={
    rectangle, draw=codeblue, line width=0.9pt,
    fill=layerA, font=\ttfamily\small,
    rounded corners=3pt, minimum width=4.4cm,
    minimum height=1.0cm, align=center, text width=4.2cm
  },
  external/.style={
    rectangle, draw=gray!60, line width=0.8pt,
    fill=gray!10, font=\small,
    rounded corners=3pt, minimum width=3.6cm,
    minimum height=0.9cm, align=center
  },
  arr/.style={-Stealth, thick, codeblue!75},
  arrgray/.style={-Stealth, thick, gray},
  reflink/.style={-Stealth, thin, codeblue!60, dashed}
}

% ── Code listings style ───────────────────────────────────────────────────────
\lstdefinestyle{generic}{
    backgroundcolor=\color{backcolour},
    commentstyle=\color{codegreen},
    keywordstyle=\color{codeblue}\bfseries,
    numberstyle=\tiny\color{codegray},
    stringstyle=\color{codegreen},
    basicstyle=\ttfamily\footnotesize,
    breakatwhitespace=false,
    breaklines=true,
    numbers=left,
    numbersep=5pt,
    showspaces=false,
    showstringspaces=false,
    tabsize=2,
    frame=single,
    rulecolor=\color{codegray},
    extendedchars=true,
    inputencoding=utf8,
    literate=%
      {á}{{\'a}}1 {č}{{\v{c}}}1 {ď}{{\v{d}}}1 {é}{{\'e}}1
      {ě}{{\v{e}}}1 {í}{{\'\i}}1 {ň}{{\v{n}}}1 {ó}{{\'o}}1
      {ř}{{\v{r}}}1 {š}{{\v{s}}}1 {ť}{{\v{t}}}1 {ú}{{\'u}}1
      {ů}{{\r{u}}}1 {ý}{{\'y}}1 {ž}{{\v{z}}}1
      {Á}{{\'A}}1 {Č}{{\v{C}}}1 {Ď}{{\v{D}}}1 {É}{{\'E}}1
      {Ě}{{\v{E}}}1 {Í}{{\'I}}1 {Ň}{{\v{N}}}1 {Ó}{{\'O}}1
      {Ř}{{\v{R}}}1 {Š}{{\v{S}}}1 {Ť}{{\v{T}}}1 {Ú}{{\'U}}1
      {Ů}{{\r{U}}}1 {Ý}{{\'Y}}1 {Ž}{{\v{Z}}}1
      {„}{{\quotedblbase}}1 {"}{{\textquotedblleft}}1
      {—}{{---}}1 {–}{{--}}1
}
\lstset{style=generic}

% ── Custom languages not bundled with listings ───────────────────────────────
\lstdefinelanguage{json}{
    keywords={true,false,null},
    keywordstyle=\color{codeblue}\bfseries,
    sensitive=true,
    string=[s]{"}{"},
    stringstyle=\color{codegreen},
    comment=[l]{//},
    commentstyle=\color{codegray},
    morestring=[b]"
}
\lstdefinelanguage{yaml}{
    keywords={name,on,push,pull_request,branches,jobs,runs-on,steps,uses,with,run},
    keywordstyle=\color{codeblue}\bfseries,
    sensitive=true,
    comment=[l]{\#},
    commentstyle=\color{codegray},
    morestring=[b]",
    morestring=[b]'
}
\lstdefinelanguage{Dockerfile}{
    keywords={FROM,RUN,COPY,WORKDIR,ENV,EXPOSE,ENTRYPOINT,CMD,LABEL,ARG,USER,VOLUME,AS},
    keywordstyle=\color{codeblue}\bfseries,
    sensitive=true,
    comment=[l]{\#},
    commentstyle=\color{codegray},
    morestring=[b]"
}

% ── Header / footer ───────────────────────────────────────────────────────────
\pagestyle{fancy}
\fancyhf{}
\fancyhead[L]{\small Petri Net Analyzer}
\fancyhead[C]{\small Závěrečná dokumentace}
\fancyhead[R]{\small Verze 1.0}
\fancyfoot[C]{\thepage}
\renewcommand{\headrulewidth}{0.4pt}

\hypersetup{
    colorlinks=true,
    linkcolor=codeblue,
    urlcolor=codeblue,
    citecolor=codeblue,
    pdftitle={Závěrečná dokumentace -- Petri Net Analyzer},
    pdfsubject={Architektura, testování a konfigurační management},
}

% ── Tolerant line-breaking + URL-style breaking inside \texttt{} ─────────────
\emergencystretch=3em
\tolerance=2000
\hyphenpenalty=2000
\exhyphenpenalty=100
% Allow line breaks at . / and other punctuation inside \texttt
\makeatletter
\def\BreakableChars{.-/_}
\makeatother

% =============================================================================
\begin{document}

% ── Title page ────────────────────────────────────────────────────────────────
\begin{titlepage}
    \centering
    \vspace*{2cm}
    {\Huge\bfseries Závěrečná dokumentace\par}
    \vspace{0.5cm}
    {\LARGE Petri Net Analyzer\par}
    \vspace{0.3cm}
    {\large Architektura, testování a konfigurační management\par}
    \vspace{1.5cm}
    {\large Verze 1.0\par}
    \vspace{0.5cm}
    {\large \today\par}
    \vfill
    \begin{tabular}{ll}
        \toprule
        \textbf{Položka} & \textbf{Hodnota} \\
        \midrule
        Projekt      & Petri Net Analyzer \\
        Dokument     & Závěrečná dokumentace \\
        Verze        & 1.0 \\
        Autor        & David Šigut \\
        Repozitář    & \url{https://github.com/seracxd/PetriNet} \\
        \bottomrule
    \end{tabular}
\end{titlepage}

\tableofcontents
\newpage

% =============================================================================
\section{Úvod}

Tento dokument popisuje výsledný stav projektu \textbf{Petri Net Analyzer} —
interaktivního editoru, simulátoru a formálního analyzátoru Petriho sítí
provozovaného jako webová aplikace. Cílem dokumentace je poskytnout úplný
pohled na tři hlavní oblasti projektu:

\begin{itemize}
    \item \textbf{Architekturu aplikace} (kapitola~\ref{sec:arch}) — zvolený
          architektonický styl, schéma komponent a použité technologie.
    \item \textbf{Testování} (kapitola~\ref{sec:tests}) — typy testů,
          použité nástroje, pokrytí kódu a strategii ověřování.
    \item \textbf{Konfigurační management} (kapitola~\ref{sec:cm}) — práce
          s verzovacím systémem, build systém, konfigurace aplikace
          a CI/CD.
\end{itemize}

Detailní popis návrhu (4+1 view, sekvenční diagramy, modely tříd) je
součástí samostatného dokumentu \emph{Software Architecture Document},
který tvoří přílohu k tomuto materiálu.

% =============================================================================
\section{Architektura aplikace}
\label{sec:arch}

\subsection{Architektonický styl}

Aplikace je realizována jako \textbf{vícevrstvý klient-server systém}
s tenkým serverem. Konkrétně se jedná o kombinaci dvou stylů:

\begin{itemize}
    \item \textbf{Modulární monolit} z pohledu nasazení — celý systém
          (server, sdílené knihovny, klient) je distribuován jako jediný
          ASP.NET Core proces běžící na serveru. Z hlediska modularity
          je rozdělen do pěti logicky oddělených .NET projektů, jejichž
          závislosti tvoří orientovaný acyklický graf.
    \item \textbf{Klient-server s tlustým klientem} z pohledu běhu —
          klientská aplikace je doručena prohlížeči jako Blazor
          WebAssembly modul a veškerá interaktivní práce (editace
          diagramu, simulace, undo/redo) probíhá lokálně v prohlížeči.
          Server slouží primárně k\,1)~prvotnímu doručení statických
          webových assetů a\,2)~výpočetně náročné formální analýze
          přenášené přes \texttt{SignalR}.
\end{itemize}

Volba tohoto stylu vychází z požadavku na minimalizaci síťové komunikace
během běžné editace (offline-friendly), ale zároveň možnosti odlehčit
slabšímu klientskému zařízení od těžkých výpočtů (analýza stavového
prostoru sítí s desítkami tisíc stavů). Pro tento dvojí provoz je
využita technologie \textbf{Blazor Auto} (interactive WebAssembly +
interactive Server render mode).

\subsection{Schéma architektury}

Architektura je rozdělena do pěti samostatných .NET projektů
(viz solution file \texttt{PetriNet.slnx}). Vrstvení a závislosti
ukazuje obrázek~\ref{fig:components}.

\begin{figure}[ht]
\centering
\resizebox{\linewidth}{!}{%
\begin{tikzpicture}[
  every node/.style={align=center},
  projectbox/.style={
    rectangle, draw=codeblue, line width=0.9pt,
    fill=layerA, font=\ttfamily\small,
    rounded corners=3pt, minimum width=4.4cm,
    minimum height=1.3cm, align=center, text width=4.2cm,
    inner sep=4pt
  },
  external/.style={
    rectangle, draw=gray!60, line width=0.8pt,
    fill=gray!10, font=\small,
    rounded corners=3pt, minimum width=3.2cm,
    minimum height=0.9cm, align=center, inner sep=4pt
  },
  arr/.style={-Stealth, thick, codeblue!75},
  arrgray/.style={-Stealth, thick, gray},
  reflink/.style={-Stealth, thin, codeblue!60, dashed}
]

  % ── Top: Browser
  \node[external, fill=layerC] (browser) at (0, 7.0)
        {Webový prohlížeč\\\scriptsize (Chrome, Firefox, Edge)};

  % ── Layer 1: Client + Server
  \node[projectbox, fill=layerA] (client) at (-4.0, 4.6)
        {PetriEditor.Client\\[1pt]\scriptsize\rmfamily Blazor WASM\\\scriptsize\rmfamily (UI, editor, simulace)};

  \node[projectbox, fill=layerB] (server) at (4.0, 4.6)
        {PetriEditor.Server\\[1pt]\scriptsize\rmfamily ASP.NET Core\\\scriptsize\rmfamily (host, SignalR, PDF)};

  % ── Layer 2: Analysis (centred)
  \node[projectbox, fill=tealpale] (analysis) at (0, 2.0)
        {PetriEditor.Analysis\\[1pt]\scriptsize\rmfamily algoritmy: stavový prostor,\\\scriptsize\rmfamily invarianty, cykly, klasifikace};

  % ── Layer 3: Shared (bottom)
  \node[projectbox, fill=layerC] (shared) at (0, -0.4)
        {PetriEditor.Shared\\[1pt]\scriptsize\rmfamily DTOs, kontrakty,\\\scriptsize\rmfamily doménové modely};

  % ── Tests (left side, separated)
  \node[projectbox, fill=layerD] (tests) at (-7.5, 2.0)
        {PetriEditor.Tests\\[1pt]\scriptsize\rmfamily xUnit\\\scriptsize\rmfamily (196 testů)};

  % ── External libraries (well separated)
  \node[external] (zblazor)  at (-8.0, 4.6) {Z.Blazor.\\Diagrams 3.0.4};
  \node[external] (signalr)  at ( 8.5, 5.6) {SignalR\\(WebSocket)};
  \node[external] (questpdf) at ( 8.5, 3.6) {QuestPDF};

  % ── Edges (browser)
  \draw[arr] (browser.south) -| node[pos=0.25,above,font=\scriptsize]{WASM}    (client.north);
  \draw[arr] (browser.south) -| node[pos=0.25,above,font=\scriptsize]{HTTPS}   (server.north);

  % ── Client <-> Server
  \draw[arr,<->] (client.east) -- node[above,font=\scriptsize]{SignalR\,/\,JSON} (server.west);

  % ── Layer 1 -> Layer 2
  \draw[arr] (client.south) -- (analysis.north west);
  \draw[arr] (server.south) -- (analysis.north east);

  % ── Layer 2 -> Shared (direct)
  \draw[arr] (analysis.south) -- (shared.north);

  % ── Client/Server bypass to Shared
  \draw[arr] (client.south) .. controls (-3.5, 1.0) and (-2.0, 0.0) .. (shared.north west);
  \draw[arr] (server.south) .. controls ( 3.5, 1.0) and ( 2.0, 0.0) .. (shared.north east);

  % ── Tests references (dashed)
  \draw[reflink] (tests.east) -- (analysis.west);
  \draw[reflink] (tests.south) .. controls (-7.5, -0.4) and (-4.0, -0.4) .. (shared.west);
  \draw[reflink] (tests.north) .. controls (-7.5,  4.6) and (-6.5,  4.6) .. (client.west);

  % ── External libs
  \draw[arrgray] (zblazor.east)  -- (client.west);
  \draw[arrgray] (signalr.west)  .. controls (6.5, 5.6) .. (server.north east);
  \draw[arrgray] (questpdf.west) -- (server.east);

\end{tikzpicture}}
\caption{Diagram komponent — pět .NET projektů a jejich vztahy
         k externím knihovnám a běhovému prostředí.
         Plné šipky = závislost na sestavení (\texttt{ProjectReference});
         čárkované šipky = reference pouze v testech.}
\label{fig:components}
\end{figure}

\paragraph{Popis projektů}
\begin{description}
    \item[PetriEditor.Shared] Sdílené doménové modely Petriho sítě
          (\texttt{Place}, \texttt{Transition}, \texttt{Arc},
          \texttt{ArcType}), DTO objekty pro SignalR přenos a mapovací
          rozhraní. Je referenční bod, na který se odkazují všechny
          ostatní projekty — tvoří \emph{spodní vrstvu} grafu závislostí.

    \item[PetriEditor.Analysis] Knihovna formální analýzy bez závislosti
          na UI nebo webovém kontextu. Obsahuje pět nezávislých analytických
          enginů (stavový prostor / Karp-Miller, P/T-invarianty, klasifikace
          struktury, Johnsonův algoritmus pro cykly, traps \& siphons)
          a sedm property testů (liveness, boundedness, safety,
          conservativeness, repetitiveness, deadlock-freedom, reversibility).
          Lze ji spustit jak v klientovi (in-process), tak na serveru.

    \item[PetriEditor.Client] Vlastní webová aplikace v Blazor WebAssembly.
          Vykresluje UI editoru (komponenty \texttt{Home.razor},
          \texttt{AnalysisPanel.razor}, \texttt{GraphTreePanel.razor}),
          spravuje stav diagramu (\texttt{PetriNetManager}), implementuje
          undo/redo (\texttt{UndoRedoService}, command pattern),
          interaktivní simulaci (\texttt{SimulationService}) a JS interop
          s knihovnou Z.Blazor.Diagrams.

    \item[PetriEditor.Server] ASP.NET Core 10 host, který doručuje
          klientské WASM assety, vystavuje \texttt{SignalR} hub
          \texttt{/hubs/analysis} pro server-side analýzu velkých sítí
          a generuje PDF exporty výsledků pomocí knihovny QuestPDF.
          Také obsahuje souborový logger a perzistenci klíčů
          \texttt{DataProtection} pro stabilitu antiforgery tokenů
          mezi restarty.

    \item[PetriEditor.Tests] Unit, integrační a regresní xUnit testy
          (196 testovacích metod, 4045 řádků) — viz kapitola~\ref{sec:tests}.
\end{description}

\subsection{Schéma rozmístění (deployment)}

Aplikace má tři typické varianty nasazení (obrázek~\ref{fig:deploy}):

\begin{enumerate}
    \item \textbf{Lokální vývoj}: \texttt{dotnet run} z projektu
          \texttt{PetriEditor.Server}; klient je servován z paměti
          a hot-reload je aktivní.
    \item \textbf{Docker kontejner}: jediný image postavený podle
          \texttt{Dockerfile} (multi-stage build, viz sekce~\ref{sec:cm-docker}),
          deployovatelný na libovolný PaaS (Railway, Azure App Service, \dots).
    \item \textbf{Statické hostování (volitelné)}: vzhledem k tomu, že
          analýza je k dispozici i ve WASM podobě, lze klientskou část
          publikovat samostatně — pak jsou těžké výpočty omezeny výkonem
          prohlížeče.
\end{enumerate}

\begin{figure}[ht]
\centering
\begin{tikzpicture}[
  every node/.style={align=center},
  procbox/.style={
    rectangle, draw=teal, fill=tealpale, line width=0.7pt,
    rounded corners=3pt, minimum width=4.4cm, minimum height=0.7cm,
    font=\footnotesize, align=center, inner sep=3pt
  },
  cli/.style={
    rectangle, draw=gray!60, fill=gray!10, line width=0.8pt,
    rounded corners=4pt, minimum width=3.6cm, minimum height=1.6cm,
    align=center, font=\small, inner sep=3pt
  },
  caplbl/.style={font=\scriptsize\itshape, gray, align=center}
]

% ── Client node (left)
\node[cli] (browser) at (-4.8, 0)
      {\textbf{Klient}\\[2pt]\scriptsize Webový prohlížeč\\\scriptsize (WASM runtime)};

% ── Server "deployment node" — UML-style box with header
\node[procbox] (kestrel) at ( 3.0,  0.30) {Kestrel HTTP/HTTPS server};
\node[procbox] (signal)  at ( 3.0, -0.55) {SignalR Hub \texttt{/hubs/analysis}};

\begin{scope}[on background layer]
  \node[draw=codeblue, line width=0.9pt, fill=layerA, rounded corners=4pt,
        fit=(kestrel)(signal),
        inner xsep=8pt, inner ysep=18pt] (host) {};
\end{scope}

\node[anchor=south, font=\bfseries\small, text=codeblue]
      at (host.north) {Server (ASP.NET Core 10)};
      
% ── Connection
\draw[arr] (browser.east) -- node[above,font=\scriptsize]{HTTPS}
                              node[below,font=\scriptsize]{WebSocket}
           (host.west);

% ── Captions below
\node[caplbl] at (-4.8, -1.6) {Blazor WASM\\doručeno přes HTTPS};
\node[caplbl] at ( 3.0, -1.6) {Docker kontejner / PaaS\\port 8080};

\end{tikzpicture}
\caption{Diagram nasazení — tenký server + tlustý WASM klient,
         obousměrné spojení přes SignalR (WebSocket).}
\label{fig:deploy}
\end{figure}

\subsection{Použité technologie a nástroje}

\begin{longtable}{@{}p{4.6cm}p{1.8cm}p{7.0cm}@{}}
    \toprule
    \textbf{Technologie} & \textbf{Verze} & \textbf{Účel} \\
    \midrule
    \endhead
    .NET                     & 10.0    & Cílová platforma všech projektů. \\
    C\#                      & 12      & Hlavní programovací jazyk. \\
    ASP.NET Core             & 10.0    & Webový server, hosting, antiforgery. \\
    Blazor WebAssembly       & 10.0    & Spuštění .NET kódu v prohlížeči. \\
    Blazor Server            & 10.0    & Interaktivní render-mode pro server-side renderování. \\
    SignalR                  & 9.0.4   & Obousměrná komunikace klient–server pro analýzu. \\
    Z.Blazor.Diagrams        & 3.0.4   & Knihovna pro vykreslování diagramů a interaktivitu. \\
    QuestPDF                 & 2024.10 & Generování PDF reportů z analýzy. \\
    DataProtection           & 9.0.3   & Trvalé klíče pro antiforgery napříč restarty. \\
    xUnit                    & 2.9.3   & Framework pro jednotkové a integrační testy. \\
    coverlet.collector       & 6.0.4   & Sběr metrik pokrytí kódu. \\
    Microsoft.NET.Test.Sdk   & 17.14.1 & Testovací běhový adapter pro \texttt{dotnet test}. \\
    Cytoscape.js (CDN)       & 3.x     & Volitelné renderování coverability grafu (s SRI hashem). \\
    Docker                   & ---     & Kontejnerizovaný build a runtime. \\
    Git + GitHub             & ---     & Verzování, code review, pull requesty. \\
    \bottomrule
\end{longtable}

\paragraph{Klíčové architektonické rozhodnutí}
Sdílení knihovny \texttt{PetriEditor.Analysis} mezi klientem i serverem
umožňuje \emph{stejný} algoritmický kód použít v obou prostředích.
Pro malé sítě (ojedinělý kontext editace) se analýza spouští přímo
v prohlížeči, pro velké sítě je delegována na server přes SignalR
s chunked přenosem grafu (200 uzlů/zpráva, viz \texttt{AnalysisHub.StreamGraph}).

% =============================================================================
\section{Testování}
\label{sec:tests}

\subsection{Strategie testování}

Projekt používá \textbf{vrstvenou strategii testování} s důrazem na
formální korektnost analytických algoritmů, neboť chybný výsledek formální
analýzy by byl pro uživatele zavádějící. Strategie se řídí třemi pravidly:

\begin{enumerate}
    \item \textbf{Algoritmické jádro je 100\% testovatelné} — knihovna
          \texttt{PetriEditor.Analysis} nemá žádné UI ani webové
          závislosti, takže všechny veřejné API jsou pokryty unit testy.
    \item \textbf{Cross-engine konzistence} — jakékoli dva engine,
          jejichž výsledky se logicky musí shodovat (např. P-invariant
          $\Rightarrow$ váha musí být zachována ve všech reachable stavech),
          jsou ověřeny v \texttt{CrossEngineConsistencyTests}.
    \item \textbf{Klasické literární vzory} — známé sítě
          (producer-consumer, vzájemné vyloučení, dining philosophers,
          bounded buffer, token ring) tvoří regresní sadu, která zajistí,
          že refaktoring algoritmu neovlivní již ověřené výsledky.
\end{enumerate}

\subsection{Typy testů}

\begin{longtable}{p{4.5cm}p{8.5cm}}
    \toprule
    \textbf{Typ} & \textbf{Příklady v projektu} \\
    \midrule
    \endhead
    \textbf{Jednotkové (unit)}      & Per-engine test soubory:
        \texttt{StateSpaceTests.cs}, \texttt{InvariantAnalysisTests.cs},
        \texttt{ClassificationAnalysisTests.cs}, \texttt{CycleDetectionTests.cs},
        \texttt{CoverabilityTreeTests.cs}, \texttt{ReachabilityTreeTests.cs},
        \texttt{SimulatorTests.cs}, \texttt{TreeLayoutEngineTests.cs},
        \texttt{GraphLayoutEngineTests.cs}. \\[3pt]
    \textbf{Property-based / Theory} &
        \texttt{TheoryTests.cs} (parametrizované testy invariantních
        vlastností: nezávislost reset arc na pořadí, hranice ořezu,
        deduplikace stavů, vážené hrany, inhibitor threshold,
        N-loop cykly). xUnit \texttt{[Theory]} + \texttt{[InlineData]}. \\[3pt]
    \textbf{Integrační}             &
        \texttt{ClassicalNetsTests.cs} — celá pipeline od konstrukce
        snapshotu po výsledný \texttt{AnalysisReport} pro pět klasických
        sítí. \\[3pt]
    \textbf{Cross-engine / regresní} &
        \texttt{CrossEngineConsistencyTests.cs} — ověřuje vzájemnou
        konzistenci výsledků (P-invarianty platí v každém dosažitelném
        stavu, coverability tree vs. state space, deduplikace cyklů,
        sémantika trapů). \\[3pt]
    \textbf{Coverage gap}            &
        \texttt{CoverageGapTests.cs}, \texttt{DefiniteDeadlockTests.cs},
        \texttt{PropertyTestsTests.cs} — explicitní testy hraničních
        případů, ořezu výpočtu a chování pod cancellation tokenem. \\[3pt]
    \textbf{End-to-end (manuální)}   &
        UI-related testy nejsou automatizovány. Manuální E2E je
        prováděno během vývoje (drag-and-drop, undo/redo, JS interop)
        v prohlížečích Chrome a Firefox. \\
    \bottomrule
\end{longtable}

\subsection{Použité nástroje}

\begin{itemize}
    \item \textbf{xUnit 2.9.3} — primární testovací framework
          (\texttt{[Fact]}, \texttt{[Theory]}, \texttt{[InlineData]},
          asserce \texttt{Assert.Equal}, \texttt{Assert.Contains},
          \texttt{Assert.True}).
    \item \textbf{Microsoft.NET.Test.Sdk 17.14.1} — adaptér umožňující
          spuštění přes \texttt{dotnet test} i z Visual Studia / JetBrains
          Rider.
    \item \textbf{xunit.runner.visualstudio 3.1.4} — VS test runner.
    \item \textbf{coverlet.collector 6.0.4} — sběr code coverage metrik;
          výstup ve formátu \texttt{cobertura.xml} (kompatibilní s ReportGenerator,
          Coveralls, Codecov).
    \item \textbf{Helpers/} — interní knihovna pro tvorbu testovacích
          sítí (\texttt{NetBuilder} fluent API), aby testy zůstaly čitelné.
\end{itemize}

\subsection{Pokrytí kódu}

\subsubsection{Statistika testů}

V projektu je definováno \textbf{196 testovacích metod}
(\texttt{[Fact]} + \texttt{[Theory]}), které dohromady reprezentují
\textbf{220 testovacích případů} po expanzi parametrizovaných
\texttt{[Theory]} přes \texttt{[InlineData]}. Rozložení podle
test-souborů ukazuje tabulka níže.

\begin{longtable}{@{}p{6.0cm}rr@{}}
    \toprule
    \textbf{Test soubor} & \textbf{\#\,Testů} & \textbf{Řádků} \\
    \midrule
    \endhead
    ClassicalNetsTests.cs                & 16   & 362 \\
    ClassificationAnalysisTests.cs       & 17   & 280 \\
    CoverabilityTreeTests.cs             & 12   & 221 \\
    CoverageGapTests.cs                  & 11   & 213 \\
    CrossEngineConsistencyTests.cs       & 13   & 335 \\
    CycleDetectionTests.cs               & 12   & 249 \\
    DefiniteDeadlockTests.cs             & \phantom{0}9    & 201 \\
    GraphLayoutEngineTests.cs            & 11   & 160 \\
    InvariantAnalysisTests.cs            & \phantom{0}9    & 187 \\
    PropertyTestsTests.cs                & 18   & 388 \\
    ReachabilityTreeTests.cs             & 13   & 243 \\
    SimulatorTests.cs                    & 24   & 401 \\
    StateSpaceTests.cs                   & 12   & 275 \\
    TheoryTests.cs                       & \phantom{0}7    & 286 \\
    TreeLayoutEngineTests.cs             & 12   & 244 \\
    \midrule
    \textbf{Celkem}                      & \textbf{196} & \textbf{4\,045} \\
    \bottomrule
\end{longtable}

\subsubsection{Naměřené pokrytí kódu}

Pokrytí bylo změřeno spuštěním všech testů s nástrojem \texttt{coverlet}
(\texttt{dotnet test --collect:"XPlat Code Coverage"}) a~následným
zpracováním \texttt{cobertura.xml} pomocí \texttt{ReportGenerator}.
Souhrnné metriky:

\begin{longtable}{@{}p{6.0cm}rr@{}}
    \toprule
    \textbf{Metrika} & \textbf{Hodnota} & \textbf{Pokryto / celkem} \\
    \midrule
    \endhead
    Line coverage    & \textbf{18.1\,\%} & 1\,485\,/\,8\,199   \\
    Branch coverage  & \textbf{15.5\,\%} & \phantom{0}711\,/\,4\,559 \\
    Method coverage  & \textbf{18.3\,\%} & \phantom{0}205\,/\,1\,119 \\
    Třídy v reportu  & 157               & --- \\
    Soubory zdrojáků & \phantom{0}72     & --- \\
    \bottomrule
\end{longtable}

Globální čísla zachycují celý solution včetně UI vrstvy a kontrakt
DTO objektů, které nemají žádnou behavior-logiku k testování.
Po rozčlenění podle assembly je pohled mnohem informativnější:

\begin{longtable}{@{}p{5.5cm}rl@{}}
    \toprule
    \textbf{Assembly} & \textbf{Line cov.} & \textbf{Účel} \\
    \midrule
    \endhead
    PetriEditor.Analysis  & \textbf{81.3\,\%} & algoritmické jádro (cíl testů) \\
    PetriEditor.Shared    & 26.0\,\%          & doménové modely + DTO kontrakty \\
    PetriEditor.Client    & \phantom{0}1.7\,\% & UI / Blazor komponenty \\
    \bottomrule
\end{longtable}

\textbf{Interpretace:}

\begin{itemize}
    \item \textbf{PetriEditor.Analysis (81.3\,\%)} — \emph{cíl
          testovacího úsilí}. Algoritmické jádro je systematicky pokryto
          jednotkovými, theory a regresními testy. Většina jednotlivých
          tříd dosahuje 90--100\,\% (viz tabulka per-class níže).
          Nepokryté zbylé části jsou:
          \begin{itemize}
              \item \texttt{Analysis.AnalysisService} (0\,\%) a~%
                    \texttt{Analysis.AnalysisReport} (0\,\%) --
                    fasáda volaná z UI; její business-logika je v~%
                    enginech, které jsou plně pokryty.
              \item \texttt{GraphLayoutBuilder} (0\,\%) -- starý kód
                    nahrazený \texttt{GraphLayoutEngine} (98.9\,\%),
                    ponechán jako fallback.
          \end{itemize}
    \item \textbf{PetriEditor.Shared (26\,\%)} — pokryty jsou doménové
          modely využívané analýzou (\texttt{PnPlace},
          \texttt{PnTransition}, \texttt{PnArc}, \texttt{PetriNetSnapshot} --
          83--100\,\%); nepokryty jsou DTO třídy
          (\texttt{*Dto.cs}) -- záměrně, jelikož jde o~čistá data-holder
          POCO bez logiky (pouze auto-generované gettery a~settery).
    \item \textbf{PetriEditor.Client (1.7\,\%)} — Blazor UI vrstva.
          Pokrytí je nízké, protože:
          \begin{itemize}
              \item Blazor komponenty (\texttt{*.razor}) vyžadují
                    rendering host (\texttt{bUnit}), který v~projektu
                    není zaveden.
              \item Drag-and-drop, JS interop a undo/redo workflow se
                    testují manuálně E2E (viz strategie testování).
              \item Jediná netriviální pokrytá třída je
                    \texttt{TreeLayoutEngine} (99\,\%), což je čistý
                    algoritmus bez UI závislosti.
          \end{itemize}
\end{itemize}

\subsubsection{Pokrytí jednotlivých tříd v PetriEditor.Analysis}

Detailní rozpis pokrytí jednotlivých tříd analytického jádra:

\begin{longtable}{@{}p{8.5cm}r@{}}
    \toprule
    \textbf{Třída} & \textbf{Line cov.} \\
    \midrule
    \endhead
    \texttt{Analysis.Algorithms.BoundednessTest}        & 100\,\% \\
    \texttt{Analysis.Algorithms.ConservativenessTest}   & 100\,\% \\
    \texttt{Analysis.Algorithms.DeadlockFreeTest}       & 100\,\% \\
    \texttt{Analysis.Algorithms.RepetitivenessTest}     & 100\,\% \\
    \texttt{Analysis.Algorithms.ReversibilityTest}      & 100\,\% \\
    \texttt{Analysis.Algorithms.SafetyTest}             & 100\,\% \\
    \texttt{Analysis.Algorithms.LivenessTest}           &  96.6\,\% \\
    \texttt{Analysis.Algorithms.FireUtils}              &  97.8\,\% \\
    \texttt{Analysis.Algorithms.GraphLayoutEngine}      &  98.9\,\% \\
    \texttt{Analysis.Algorithms.CoverabilityTreeBuilder}&  91.1\,\% \\
    \texttt{Analysis.Algorithms.ReachabilityTreeBuilder}&  94.7\,\% \\
    \texttt{Analysis.Engines.ClassificationAnalysis}    &  96.9\,\% \\
    \texttt{Analysis.Engines.CyclesAnalysis}            &  95.5\,\% \\
    \texttt{Analysis.Engines.InvariantAnalysis}         &  87.6\,\% \\
    \texttt{Analysis.Engines.StateSpaceAnalysis}        &  96.0\,\% \\
    \texttt{Analysis.Engines.TrapCotrapAnalysis}        & 100\,\% \\
    \texttt{Analysis.Simulation.PetriNetSimulator}      & 100\,\% \\
    \texttt{Analysis.AnalysisBundle}                    & 100\,\% \\
    \bottomrule
\end{longtable}

Plný HTML report (s navigací do zdrojového kódu, zvýrazněnými
nepokrytými řádky a~grafy historického trendu) je vygenerován do
\texttt{docs/coverage/index.html}.

\subsubsection{Reprodukce výsledků}

Lokální spuštění všech testů s metrikou pokrytí:

\begin{lstlisting}[language=bash]
# Vsechny testy
dotnet test

# Vcetne sberu pokryti (coverlet)
dotnet test --collect:"XPlat Code Coverage"

# Volitelne: HTML report z cobertura.xml
dotnet tool install -g dotnet-reportgenerator-globaltool
reportgenerator -reports:"**/coverage.cobertura.xml" \
                -targetdir:"coveragereport" \
                -reporttypes:Html
\end{lstlisting}

\textbf{Výsledek běhové statistiky:} 220/220 testů prochází za $\sim$2~s
(release build, Windows 11, .NET 10 SDK). Build je čistý — 0 chyb,
1 varování (NU1510 package-prune, neblokující).

\subsection{Ukázky testů}

\subsubsection{Jednotkový test — invariantní vlastnost simulátoru}

Test ověřuje, že \texttt{PetriNetSimulator} po vyresetování vrátí
síť do iniciálního stavu nezávisle na předchozí historii vystřelení:

\begin{lstlisting}[language={[Sharp]C}]
[Fact]
public void Reset_RestoresInitialMarking()
{
    var net = NetBuilder.Linear(p1: 3, p2: 0);
    var sim = new PetriNetSimulator(net);

    sim.FireAll();             // vyprazdni p1
    Assert.Equal(0, sim.Tokens["p1"]);

    sim.Reset();
    Assert.Equal(3, sim.Tokens["p1"]);
    Assert.Empty(sim.History);
}
\end{lstlisting}

\subsubsection{Theory test — sémantika reset hran}

Parametrizovaný test ověřuje, že pořadí výpisu hran (normal vs. reset)
nemá vliv na výsledný marking:

\begin{lstlisting}[language={[Sharp]C}]
[Theory]
[InlineData(ArcOrder.NormalThenReset)]
[InlineData(ArcOrder.ResetThenNormal)]
public void ResetArc_OrderIndependent(ArcOrder order)
{
    var net = NetBuilder.MixedArcs(order, initialTokens: 5);
    var sim = new PetriNetSimulator(net);

    sim.Fire("t1");

    Assert.Equal(0, sim.Tokens["p1"]);   // reset vyhrava
    Assert.Equal(1, sim.Tokens["p2"]);
}
\end{lstlisting}

\subsubsection{Cross-engine konzistence}

Test ověřuje, že každý spočítaný P-invariant skutečně platí ve všech
dosažitelných stavech (definice P-invariantu: $y \cdot M = y \cdot M_0$):

\begin{lstlisting}[language={[Sharp]C}]
[Fact]
public void PInvariants_HoldAtEveryReachableMarking()
{
    var net = NetBuilder.MutualExclusion();
    var bundle = AnalysisService.Run(net.ToSnapshot());

    foreach (var pi in bundle.Invariants.PInvariants)
    foreach (var marking in bundle.StateSpace.Markings)
    {
        var weighted = DotProduct(pi, marking);
        var initial  = DotProduct(pi, bundle.StateSpace.Initial);
        Assert.Equal(initial, weighted);
    }
}
\end{lstlisting}

\subsection{Strategie testování — kdy a jak}

\begin{itemize}
    \item \textbf{Při vývoji:} test-driven přístup pro čisté algoritmy
          (state space, invariants); nejprve napsán test podle
          definice z literatury (Murata 1989), poté implementace.
    \item \textbf{Před commit:} \texttt{dotnet test} musí projít s 0 selháními;
          \texttt{dotnet build /p:TreatWarningsAsErrors=true} se používá
          jako lint.
    \item \textbf{Při bug fixu:} je-li to možné, nejprve přidán failing
          test reprodukující bug, teprve pak oprava (regresní test).
          Příklad: \texttt{TheoryTests.cs} obsahuje regresní test pro
          dříve špatné chování reset hran v závislosti na pořadí výpisu.
    \item \textbf{Při refaktoringu velkých částí:} cross-engine testy
          jsou primární bezpečnostní síť; pokud byly před refaktorem
          zelené a zůstávají zelené po refaktoru, lze refaktor považovat
          za sémanticky neutrální.
    \item \textbf{Manuální E2E testování:} prováděno v prohlížeči
          podle scénářů popsaných v Software Architecture Document
          (drag-and-drop arc, undo/redo deep stack, simulace s
          time-travel, dock/float panelu, JS interop pro klávesové
          zkratky).
\end{itemize}

% =============================================================================
\section{Konfigurační management}
\label{sec:cm}

\subsection{Verzovací systém}

Projekt používá \textbf{Git} jako verzovací systém s repozitářem
hostovaným na GitHubu (\url{https://github.com/seracxd/PetriNet}).

\subsubsection{Struktura repozitáře}

\begin{lstlisting}[language=bash]
PetriNet/
|-- .github/
|   `-- workflows/              # Pripraveno pro CI (zatim prazdne)
|-- .gitattributes              # Normalizace EOL a binarnich souboru
|-- .gitignore                  # Standardni VS / dotnet ignore
|-- Dockerfile                  # Multi-stage build & deploy
|-- PetriNet.slnx               # SLN replacement (novy .slnx format)
|-- TODO.md                     # Roadmap, zmeny, code-review log
|-- PetriEditor.Shared/         # Domenove modely a kontrakty
|-- PetriEditor.Analysis/       # Algoritmy formalni analyzy
|-- PetriEditor.Client/         # Blazor WASM klient
|-- PetriEditor.Server/         # ASP.NET Core host
`-- PetriEditor.Tests/          # xUnit testy (196)
\end{lstlisting}

\subsubsection{Větvení a workflow}

Aplikuje se zjednodušený \textbf{trunk-based} model:

\begin{itemize}
    \item \texttt{master} — hlavní větev, vždy buildovatelná, vždy
          s procházejícími testy (32 commitů celkem v čase psaní dokumentu).
    \item \texttt{feature/*} — pracovní větve pro nové funkce
          (např. \texttt{feature/manual-graph-layout} obsahuje refaktor
          grafu z Cytoscape na vlastní canvas renderer).
    \item \texttt{archive/*} — archivované předchozí varianty (např.
          \texttt{archive/original-blazor-wasm} obsahuje původní
          single-project verzi před modularizací).
\end{itemize}

Slévání do \texttt{master} probíhá výhradně přes \textbf{pull request}
(\#1, \#2 v historii) s manuální revizí změn. Velké refaktory jsou
prováděny v sérii malých commitů s popisnými zprávami:

\begin{lstlisting}[language=bash]
2fc6285  Refactor graph view: new layout, canvas, remove Cytoscape
cea84d9  Hide uncertain analyses for nets with inhibitor/reset arcs
c9e88d1  Fix DI scope, timer disposal, nullability, and error logging
8777e1d  Auto-layout, chunked graph streaming, and test coverage
4d726b7  Add integration and regression tests for analysis engines
\end{lstlisting}

\subsection{Build systém}

Projekt používá nativní \textbf{.NET CLI} (\texttt{dotnet}) jako
build systém. Žádné externí orchestrátory (Make, Cake, NUKE)
se nepoužívají — výhodou je nulová nárazová křivka pro každého
.NET vývojáře a přímá integrace s IDE.

\subsubsection{Klíčové příkazy}

\begin{longtable}{@{}p{7.6cm}p{5.9cm}@{}}
    \toprule
    \textbf{Příkaz} & \textbf{Účel} \\
    \midrule
    \endhead
    \texttt{dotnet restore}                                            & Stažení NuGet balíčků pro celé řešení. \\
    \texttt{dotnet build}                                              & Sestavení všech projektů (Debug). \\
    \texttt{dotnet build -c Release}                                   & Optimalizovaný release build. \\
    \texttt{dotnet build /p:TreatWarningsAsErrors=true}                & „Lint“ režim — varování blokují build. \\
    \texttt{dotnet test}                                               & Spuštění 220 xUnit test. případů. \\
    \texttt{dotnet test -{-}collect:'XPlat Code Coverage'}             & Testy + sběr code coverage. \\
    \texttt{dotnet run -{-}project PetriEditor.Server}                 & Lokální spuštění (dev. server). \\
    \texttt{dotnet publish -c Release -o publish/}                     & Publikace pro nasazení (self-contained výstup). \\
    \bottomrule
\end{longtable}

\subsubsection{NuGet závislosti}

Externí balíčky jsou definovány v \texttt{.csproj} souborech
(jeden centralní \texttt{Directory.Packages.props} se nepoužívá kvůli
malému počtu závislostí). Klíčové balíčky:

\begin{longtable}{@{}>{\raggedright\arraybackslash}p{7.6cm}p{1.6cm}p{4.5cm}@{}}
    \toprule
    \textbf{Balíček} & \textbf{Verze} & \textbf{Použito v} \\
    \midrule
    \endhead
    Microsoft.AspNetCore.Components.WebAssembly        & 10.0.2     & Client \\
    Microsoft.AspNetCore.Components.WebAssembly.Server & 10.0.2     & Server \\
    Microsoft.AspNetCore.SignalR.Client                & 9.0.4      & Client \\
    Microsoft.AspNetCore.DataProtection                & 9.0.3      & Server \\
    QuestPDF                                           & 2024.10.4  & Server (PDF export) \\
    Z.Blazor.Diagrams (+ Core)                         & 3.0.4      & Client \\
    xunit                                              & 2.9.3      & Tests \\
    coverlet.collector                                 & 6.0.4      & Tests \\
    Microsoft.NET.Test.Sdk                             & 17.14.1    & Tests \\
    \bottomrule
\end{longtable}

\subsection{Konfigurace aplikace}

Aplikace pracuje s konfigurací na třech úrovních:

\subsubsection{appsettings.json (server)}

Serverový projekt používá standardní ASP.NET Core konfigurační hierarchii:
\texttt{appsettings.json} (defaultní) $\to$ \texttt{appsettings.Development.json}
(override) $\to$ environment variables (nejvyšší priorita).

\begin{lstlisting}[language=json]
{
  "Logging": {
    "LogLevel": {
      "Default": "Information",
      "Microsoft.AspNetCore": "Warning"
    }
  },
  "AllowedHosts": "*"
}
\end{lstlisting}

\subsubsection{Environment variables}

Pro nasazení na PaaS s ephemerálním souborovým systémem (Railway, Heroku)
program rozpoznává tyto proměnné prostředí:

\begin{longtable}{p{4.5cm}p{8.5cm}}
    \toprule
    \textbf{Proměnná} & \textbf{Účel} \\
    \midrule
    \endhead
    \texttt{ASPNETCORE\_ENVIRONMENT}    & \texttt{Development} / \texttt{Production} (řídí WASM debugger, exception handler). \\
    \texttt{DP\_KEYS\_PATH}             & Cesta pro perzistenci DataProtection klíčů (default: \texttt{AppContext.BaseDirectory/dataprotection-keys}). \\
    \texttt{LOW\_MEMORY\_GC}            & Zapne \texttt{GCLatencyMode.Batch} pro hostingy s omezenou pamětí. \\
    \texttt{RAILWAY\_ENVIRONMENT}       & Auto-detekce Railway hostingu (alias pro \texttt{LOW\_MEMORY\_GC}). \\
    \bottomrule
\end{longtable}

\subsubsection{Klientská konfigurace}

Klient (Blazor WASM) má oddělený konfigurační kanál:

\begin{itemize}
    \item \texttt{DiagramSettings} — runtime singleton držící vizuální
          a behaviorální parametry (velikost grid, barvy hran, velikosti
          uzlů). Je inicializován z \texttt{appsettings.json} přes
          \texttt{IConfiguration.Bind} a má event \texttt{Changed}
          pro re-render.
    \item \texttt{localStorage} (browser) — perzistentní uživatelské
          preference (firing delay simulace, zapnutí/vypnutí grid
          overlay). Není to server-side state — žádný cookies,
          žádná databáze.
\end{itemize}

\subsection{Dockerizace}
\label{sec:cm-docker}

Projekt je distribuovatelný jako jediný Docker image. \texttt{Dockerfile}
používá \textbf{multi-stage build}, aby výsledný image obsahoval pouze
runtime artefakty bez SDK:

\begin{lstlisting}[language=Dockerfile]
FROM mcr.microsoft.com/dotnet/sdk:10.0 AS build
WORKDIR /src
COPY . .
RUN dotnet test PetriEditor.Tests/PetriEditor.Tests.csproj \
        -c Release --no-restore
RUN dotnet publish PetriEditor.Server/PetriEditor.Server.csproj \
        -c Release -o /app/publish

FROM mcr.microsoft.com/dotnet/aspnet:10.0 AS runtime
WORKDIR /app
COPY --from=build /app/publish .
ENV ASPNETCORE_ENVIRONMENT=Production
EXPOSE 8080
ENTRYPOINT ["dotnet", "PetriEditor.Server.dll"]
\end{lstlisting}

\textbf{Klíčové body:}
\begin{itemize}
    \item \textbf{Testy jsou součástí buildu} (řádek~4): pokud
          \texttt{dotnet test} selže, image se nepostaví — bránění
          chybnému deploy.
    \item \textbf{Two-stage build}: výsledný image (\texttt{aspnet:10.0})
          je menší než SDK image (\texttt{sdk:10.0}). Reálná velikost
          po publikaci $\sim$\,200~MB.
    \item \textbf{Production environment} fixed (řádek~10) — žádný
          accidentální development režim.
\end{itemize}

\subsection{CI/CD}

V současné fázi projekt \textbf{nepoužívá automatizovanou CI/CD pipeline}.
Adresář \texttt{.github/workflows/} je vytvořen, ale prázdný. Důvody:

\begin{itemize}
    \item Projekt je jednočlenný (akademický kontext), takže nedochází
          k souběžnému psaní více vývojáři, kde CI brání rozbití
          \texttt{master}u.
    \item \textbf{Dockerfile sám o sobě supluje CI} — testy jsou
          součástí buildu (řádek 4 \texttt{Dockerfile}), takže každé
          provedené nasazení implicitně garantuje zelené testy.
    \item Lokální workflow je jednoduchý:
          \texttt{dotnet test} $\to$ \texttt{git commit} $\to$
          \texttt{git push} $\to$ PR review $\to$ merge.
\end{itemize}

\paragraph{Plán pro CI/CD}
Do připraveného adresáře \texttt{.github/workflows/} lze umístit
GitHub Actions soubor \texttt{ci.yml} přibližně v této podobě:

\begin{lstlisting}[language=yaml]
name: CI
on:
  push:
    branches: [master]
  pull_request:
    branches: [master]
jobs:
  build-test:
    runs-on: ubuntu-latest
    steps:
      - uses: actions/checkout@v4
      - uses: actions/setup-dotnet@v4
        with:
          dotnet-version: 10.0.x
      - run: dotnet restore
      - run: dotnet build --no-restore -c Release
      - run: dotnet test --no-build -c Release \
               --collect:"XPlat Code Coverage"
      - uses: codecov/codecov-action@v4
\end{lstlisting}

Tento workflow zajistí, že každý PR proti \texttt{master} bude před
slevením automaticky ověřen.

% =============================================================================
\section{Bezpečnostní opatření a omezení}
\label{sec:security}

Přestože aplikace neukládá uživatelská data ani neimplementuje
autentizaci (jde o bezstavovou nástrojovou aplikaci), na serverové
straně jsou implementována standardní opatření proti zneužití
veřejně exponovaného endpointu:

\begin{description}
    \item[Antiforgery + DataProtection] ASP.NET Core antiforgery
          tokeny jsou aktivní (\texttt{app.UseAntiforgery()}).
          Klíče \texttt{DataProtection} jsou perzistovány na disk
          (\texttt{DP\_KEYS\_PATH}), takže tokeny vystavené před
          restartem serveru zůstávají platné pro již aktivní
          klientské sessions.

    \item[Limit velikosti vstupní sítě] \texttt{AnalysisHub.ValidateNetSize}
          odmítne sítě s více než $100\,000$ prvky (\texttt{MaxNetElements})
          výjimkou \texttt{HubException} ještě před spuštěním analýzy.
          Brání DoS přes obří JSON payload.

    \item[Globální semafor + per-connection throttling]
          \texttt{SemaphoreSlim} velikosti $\max(2, \text{ProcessorCount})$
          omezuje souběžný počet těžkých analytických operací na celém
          serveru. Na úrovni jednotlivého SignalR connection je nastaven
          250\,ms throttle, takže jeden klient nemůže zaplavit hub
          opakovanými voláními.

    \item[Wall-clock deadline 60\,s] Každá hub-metoda
          (\texttt{RunAnalysis}, \texttt{RunGraphAnalysis},
          \texttt{ExportPdf}) má vlastní \texttt{CancellationTokenSource}
          s deadline 60\,s. Po jeho vypršení je celý výpočet zrušen
          a klientovi se vrátí \texttt{HubException}.

    \item[SignalR receive size cap] \texttt{MaximumReceiveMessageSize}
          je nastaveno na 10\,MB. Větší zprávy jsou zahozeny dříve,
          než se začnou deserializovat.

    \item[Subresource Integrity] Externí Cytoscape.js načítaný z CDN
          má \texttt{integrity="sha384-..."}, \texttt{crossorigin}
          a \texttt{referrerpolicy="no-referrer"}. Pokud by byl CDN
          obsah nahrazen, prohlížeč skript odmítne.
\end{description}

\textbf{Známá omezení (mimo bezpečnost):}

\begin{itemize}
    \item Aplikace \emph{neimplementuje persistenci} — diagramy se po
          přenačtení stránky ztratí. Plánovaná infrastrukturní vrstva
          (sekce v SAD) by tento bod adresovala.
    \item Není podpora \emph{časovaných Petriho sítí}; priority
          přechodů jsou implementovány, ale nedeterministická volba
          při shodné prioritě je řešena pevným pořadím dle ID.
    \item \emph{Coverage UI vrstvy je nízká} — Blazor komponenty
          nejsou pokryty automatizovanými testy (chybí \texttt{bUnit}
          dependency). Manuální E2E testování v Chrome a~Firefox je
          dokumentováno v sekci~\ref{sec:tests}.
\end{itemize}

% =============================================================================
\section{Závěr}

Petri Net Analyzer je implementován jako modulární Blazor Auto aplikace
rozdělená do pěti .NET projektů. Jádro aplikace (\texttt{PetriEditor.Analysis})
je čistá doménová knihovna bez vazby na UI ani web, což umožňuje její
sdílení mezi klientem (WebAssembly) a serverem (ASP.NET Core)
a zároveň úplné pokrytí jednotkovými testy.

Testovací strategie se opírá o 196 xUnit testovacích metod ve formě
220 testovacích případů (4045 řádků) rozdělených do pěti kategorií
(jednotkové, theory, integrační, cross-engine, coverage-gap) —
všechny procházejí v $\sim$2 sekundách. Pokrytí algoritmického jádra
\texttt{PetriEditor.Analysis} dosahuje \textbf{81.3\,\%} (line coverage),
přičemž většina jednotlivých enginů má 95--100\,\%. Build a deploy
je zajištěn nativním .NET CLI a multi-stage Docker imagem, který testy
exekvuje jako součást build pipeline. Verzování probíhá v Gitu na GitHubu
v trunk-based modelu s pull-request review.

Detailní popis jednotlivých návrhových rozhodnutí, sekvenčních diagramů
a class modelu je obsažen v doprovodném \emph{Software Architecture Document}u.

\end{document}
